\section{The data pipeline: representation, algorithm, and failure modes}
\label{sec:pipeline-audit}

\subsection{Evidence and scope}
This section audits the retrieved ASDMotion source at commit
\texttt{f699dc7} and its MMAction2 fork at \texttt{f632b7e}.
These are retrieved repository versions, not verified paper-era execution
versions. Source snapshots and full line-level evidence are preserved in
\texttt{references/source\_snapshots/} and
\texttt{reports/pipeline\_audit.md}. Ten executable counterexample tests
in \texttt{tests/test\_upstream\_semantics.py} pass; they reproduce the
documented upstream behaviors without CUDA or MMCV. A passing
counterexample test does not mean that the underlying behavior is correct.
No claim of reproduced training or published accuracy is made here.
The source references are the
\href{https://github.com/Dinstein-Lab/ASDMotion/tree/f699dc7b7af5726df5d8e79161d192eac22cfafb}{pinned ASDMotion repository}
and the
\href{https://github.com/TalBarami/mmaction2/tree/f632b7e56291805ac6d4debc8c7421b0d2c48f2e}{pinned MMAction2 fork}.

\subsection{From pixels to a movement score}
The inspected fresh-video path extracts BODY\_25 OpenPose coordinates and
confidence, converts them to COCO-17, and keeps up to five person slots.
Each frame independently sorts skeletons by mean confidence, so slot index
does not identify a persistent person. YOLOv5 child boxes are matched to
skeleton boxes by intersection over union, followed by nearby-frame
recovery. The fork's \texttt{ChildDetect} transform then selects one slot
per frame. These are distinct operations: pose estimation provides body
geometry, while child selection determines whose geometry is analyzed.
\footnote{ASDMotion \texttt{openpose\_executor.py:171--181};
\texttt{skeleton\_matcher.py:19--74}; fork
\texttt{pose\_loading.py:699--707}.}

The COCO-17 joint order is nose, left/right eye, left/right ear,
left/right shoulder, left/right elbow, left/right wrist, left/right hip,
left/right knee, and left/right ankle. Its indices in BODY\_25 are
\texttt{[0,16,15,18,17,5,2,6,3,7,4,12,9,13,10,14,11]}.
Thus the model receives neither explicit finger landmarks nor toe/heel,
neck, or mid-hip landmarks. Coordinates and confidence are separate arrays,
normally $X\in\mathbb{R}^{M\times T\times17\times2}$ and
$Q\in\mathbb{R}^{M\times T\times17}$. These arrays contain two-dimensional
image coordinates, not measured three-dimensional body positions.
\footnote{ASDMotion \texttt{skeleton\_layout.py:35--38,41--75,114--138};
\texttt{openpose\_executor.py:171--198}.}

The default inference windows are 200 frames with a 30-frame stride and
minimum length 60. At 30 frames per second, these correspond to about
6.67, 1, and 2 seconds; the durations change if frame rate differs.
Tail windows can be shorter. The fork samples 48 frames from each sequence,
spatially compacts and resizes poses, and renders 17 confidence-weighted
Gaussian heatmap channels with standard deviation 0.6. Training uses
$48\times48$ maps and test uses $56\times56$ maps. Limb maps are disabled.
\footnote{ASDMotion \texttt{splitter.py:8--37},
\texttt{preprocess.py:79}; fork \texttt{asdmotion.py} pipelines and
\texttt{pose\_loading.py:15--131,407--450,583--619}.}

For intuition, a joint heatmap has the form
\begin{equation}
 H_{t,j}(u,v)=\max_m q_{m,t,j}
 \exp\left[-\frac{(u-x_{m,t,j})^2+(v-y_{m,t,j})^2}{2\sigma^2}\right].
\end{equation}
Here $m$ indexes people (one after successful child selection), $t$ sampled
frames, $j$ joints, $(u,v)$ heatmap pixels, and $q$ pose-estimator confidence.
This formula summarizes the Gaussian-and-maximum implementation, which
also truncates rendering to a local neighborhood.
\footnote{Fork \texttt{pose\_loading.py:407--450}.}
The recognizer applies three-dimensional convolutions across heatmap space
and time: ResNet3dSlowOnly depth 50 with a two-class I3D head, initialized
from Kinetics-400 weights. Ten temporal samples and horizontal flips are
averaged at test time. This is a spatiotemporal classifier over 2D poses,
not a 3D pose estimator.\footnote{Fork \texttt{asdmotion.py}, model,
test pipeline, and \texttt{load\_from}.}

\paragraph{Interpretation.}
The representation makes posture, body-joint displacement and coordination
available to the classifier. It does not explicitly encode object identity,
speech, finger articulation, or the child's reason for moving. Whether
such missing information limits a proposed target is an experimental
question. The published task is SMM detection within children with ASD,
not ASD-versus-non-ASD diagnosis.\cite{barami2024}

\subsection{Clip score, event, and burden are different targets}
The paper assigns each frame the maximum score of its covering windows,
thresholds at 0.85 inclusively, then merges contiguous positive frames.
\cite{barami2024} A transparent reference evaluator should implement
\begin{equation}
 p_t=\max_{w:t\in w}p_w,\qquad
 z_t=\mathbf{1}\{p_t\geq\tau\},\qquad \tau=0.85,
\end{equation}
with an explicit uncovered-frame policy. Event count is the number of
positive runs, whereas duration is the sum of positive frame durations.
For an explicitly observed-time estimand, a proposed valid-frame burden is
\begin{equation}
 B_{\mathrm{obs}}=
 \frac{\sum_t v_t z_t\,\Delta t_t}{\sum_t v_t\,\Delta t_t},
\end{equation}
where $v_t$ is a documented validity mask and $\Delta t_t$ is frame duration.
This proposed masked estimand must be distinguished from full-recording
burden; missing observations are not automatically negative observations.
It is motivated by the source summary dividing complete event spans by
valid-child-frame count without intersecting the numerator with validity.
\footnote{ASDMotion \texttt{detector.py:70--89}.}

\subsection{Confirmed implementation findings}
The following are findings about the pinned public source, not evidence
that the published analysis used these same defects.
\begin{enumerate}
\item \textbf{Fresh-path execution blockers.} Aggregation references
undefined \texttt{NET\_NAME}. Pose conversion reads an \texttt{adjust}
key absent from the dictionary generated immediately upstream.
Disabling the advertised optional child detector sets all IDs to $-1$
and subsequently raises a no-child exception.
\footnote{\texttt{aggregator.py:62};
\texttt{openpose\_executor.py:105--113,196};
\texttt{preprocess.py:58--64}. The first two are tested counterexamples.}
\item \textbf{Event aggregation error.} After bypassing the undefined
name in the test harness, scores for positive $[0,200)$, negative
$[30,230)$, and positive $[60,260)$ yield two overlapping positive events
and an intervening negative-duration interval, rather than one union
$[0,260)$. The last positive run is excluded by the merge-loop bound.
The implementation also uses $>$ instead of the paper's $\geq$.
\footnote{\texttt{aggregator.py:20,44}; executable tests
\texttt{test\_final\_positive\_run\_not\_merged} and
\texttt{test\_threshold\_equality\_excluded}.}
\item \textbf{Missing identity is not safely represented.}
Fork \texttt{ChildDetect} indexes an ID of $-1$ as the last person slot.
Matcher recovery chooses the farther of two neighboring reference frames;
frame index zero is also treated as false in a proximity condition.
The first two behaviors have executable counterexamples.
\footnote{Fork \texttt{pose\_loading.py:705--706}; ASDMotion
\texttt{skeleton\_matcher.py:58--60}.}
\item \textbf{Burden reporting has edge cases.} A recording with no
positive event produces no summary row. A 300-frame positive interval
divided by only 100 valid child frames produces a burden of 3.0.
Both are demonstrated by tests; their frequency in real data is unknown.
\footnote{ASDMotion \texttt{detector.py:71--82}.}
\item \textbf{Long-clip evaluation can depend on random cropping.}
The fork's test pipeline applies a Python-random 200-frame crop before
sampling. The sampler fixes a NumPy seed, which does not by itself control
Python's random generator. Caller-wide seeding may control the crop, but
the config does not exhaustively evaluate a long clip.
\footnote{Fork \texttt{augmentations.py:1962--1983},
\texttt{pose\_loading.py:79}, and \texttt{asdmotion.py} test pipeline.}
\end{enumerate}

\subsection{Reproducibility risks and adapter requirements}
Additional source risks are unsorted pose-JSON file enumeration, accepted
but unreconciled detection/pose length mismatch, width--height convention
mismatch on fresh extraction, machine-specific dependency paths and conda
launcher, removed NumPy \texttt{np.int} usage, and existence-only caches.
The line-level audit separates confirmed interface differences from risks
whose occurrence in the released data is unknown.
\footnote{ASDMotion \texttt{openpose\_executor.py:126,191};
\texttt{utils.py} video properties;
\texttt{skeleton\_matcher.py:82--99}; \texttt{requirements.txt};
\texttt{resources/run\_in\_env.bat}; fork
\texttt{pose\_loading.py:127}; ASDMotion
\texttt{preprocess.py:33--54}, \texttt{detector.py:30,61--65}.}

The annotated-release schema must be adapted before using this fork.
Its loader expects \texttt{label}, \texttt{frame\_dir}/\texttt{filename},
and split-member lists referring to sample identifiers. Already-selected
child arrays of shape $[T,17,2]$ need a singleton person axis and should
bypass \texttt{ChildDetect}. Split names must match actual archive keys,
not copied config placeholders. Verify binary-label semantics, unique
sample identity, participant grouping, frame rates and time units first.
Do not interpret a field as frames solely because its name contains
\texttt{frame}; reconcile it against array length and frame rate.
\footnote{Fork \texttt{pose\_dataset.py:98--113} and
\texttt{pose\_loading.py:699--707}; proposed adapter requirements, with
released-data measurements documented separately in the data audit.}

The config's loss weights are $[0.75,0.25]$, while inference treats class
1 as SMM. If training uses the same mapping, the positive class receives
one third the per-example negative-class weight. This is not automatic
minority-class compensation. Preserve the original setting for a named
reproduction experiment, and test alternative weighting separately after
confirming the released labels. The config also uses the same
\texttt{test0} split for validation and test, so model-selection procedures
need independent clarification.\footnote{Fork \texttt{asdmotion.py};
ASDMotion \texttt{detector.py:49--50}.}

\paragraph{Recommended next step.}
Before adding a more complex model, create a validated deterministic
adapter and a separate paper-defined evaluator. Preserve explicit unknown
identity, test empty and overlapping-event cases, and establish independent
participant splits. Then compare simple kinematic features and PoseC3D
under the same coverage and duration metrics. This is a proposed research
sequence motivated by the audit, not an experimental result.
